\documentclass[aspectratio=169,11pt]{beamer}
% Shared Beamer style for practice contest solution walkthroughs.
\usepackage[utf8]{inputenc}
\usepackage[T1]{fontenc}
\usepackage{lmodern}
\usepackage{amsmath,amssymb}
\usepackage{xcolor}
\usepackage{hyperref}
\usepackage{listings}

\definecolor{CPNavy}{HTML}{172033}
\definecolor{CPBlue}{HTML}{2563EB}
\definecolor{CPGreen}{HTML}{16A34A}
\definecolor{CPOrange}{HTML}{F97316}
\definecolor{CPGray}{HTML}{64748B}
\definecolor{CPLight}{HTML}{F8FAFC}
\definecolor{CPCodeBg}{HTML}{F1F5F9}
\definecolor{CPCodeRule}{HTML}{CBD5E1}

\usetheme{default}
\usefonttheme{professionalfonts}
\setbeamertemplate{navigation symbols}{}
\setbeamersize{text margin left=0.65cm,text margin right=0.65cm}
\setbeamertemplate{itemize item}{\color{CPBlue}$\blacktriangleright$}
\setbeamertemplate{itemize subitem}{\color{CPGray}$\triangleright$}

\setbeamercolor{normal text}{fg=CPNavy,bg=white}
\setbeamercolor{frametitle}{fg=white,bg=CPNavy}
\setbeamercolor{title}{fg=white,bg=CPNavy}
\setbeamercolor{subtitle}{fg=CPLight,bg=CPNavy}
\hypersetup{colorlinks=true,urlcolor=CPBlue}

\setbeamertemplate{frametitle}{%
  \nointerlineskip%
  \begin{beamercolorbox}[wd=\paperwidth,ht=0.88cm,dp=0.22cm,leftskip=0.55cm,rightskip=0.35cm]{frametitle}
    \usebeamerfont{frametitle}\insertframetitle\hfill\small\insertframenumber/\inserttotalframenumber
  \end{beamercolorbox}%
}

\setbeamertemplate{footline}{%
  \leavevmode%
  \hbox{%
    \begin{beamercolorbox}[wd=.58\paperwidth,ht=2.4ex,dp=1ex,leftskip=.5cm]{author in head/foot}%
      \color{CPGray}\insertshorttitle
    \end{beamercolorbox}%
    \begin{beamercolorbox}[wd=.42\paperwidth,ht=2.4ex,dp=1ex,rightskip=.5cm plus1fil]{date in head/foot}%
      \hfill\color{CPGray}\insertshortauthor
    \end{beamercolorbox}%
  }%
}

\setbeamertemplate{title page}{%
  \vspace*{1.25cm}
  \begin{beamercolorbox}[wd=\paperwidth,sep=0.65cm,leftskip=0.15cm,rightskip=0.15cm]{title}
    {\color{white}\Huge\bfseries\inserttitle\par}
    \vspace{0.35cm}
    {\color{CPLight}\Large\insertsubtitle\par}
    \vspace{0.6cm}
    {\color{CPLight}\large\insertauthor\par}
  \end{beamercolorbox}
  \vfill
}

\newcommand{\sectiontag}[1]{\textbf{\color{CPBlue}#1}}
\newenvironment{tightitemize}{\begin{itemize}\small\setlength{\itemsep}{0.18cm}}{\end{itemize}}
\lstdefinestyle{cpcode}{
  language=Python,
  basicstyle=\ttfamily\tiny,
  keywordstyle=\color{CPBlue}\bfseries,
  commentstyle=\color{CPGray}\itshape,
  stringstyle=\color{CPGreen},
  backgroundcolor=\color{CPCodeBg},
  frame=single,
  framerule=0.25pt,
  rulecolor=\color{CPCodeRule},
  showstringspaces=false,
  columns=fullflexible,
  keepspaces=true,
  breaklines=true,
  breakatwhitespace=false,
  tabsize=4,
  xleftmargin=0.08cm,
  xrightmargin=0.08cm,
  aboveskip=0.08cm,
  belowskip=0.08cm
}


\title[10kindsofpeople]{10 Kinds of People}
\subtitle{Day 4 solution walkthrough}
\author[Solution Deck]{Competitive Programming Bootcamp}
\date{}

\begin{document}

\begin{frame}[plain]
  \titlepage
\end{frame}

% BEGIN GENERATED STATEMENT INTERPRETATION
\begin{frame}{Interpret the Word Problem}
\small
\sectiontag{Statement excerpts:}
\begin{tightitemize}
  \item \textit{``Users of binary have to stay in the zones marked with a zero''}
  \item \textit{``People can move north, south, east or west''}
\end{tightitemize}

\vspace{0.18cm}
\sectiontag{Programming interpretation:}
\begin{tightitemize}
  \item The grid is a graph where adjacent equal digits are connected.
  \item Precompute connected components once instead of searching for every query.
  \item Two queried cells are connected only if their component IDs match.
\end{tightitemize}
\end{frame}
% END GENERATED STATEMENT INTERPRETATION







\begin{frame}{Problem Model}
\small
\sectiontag{Textbook connection:} CP4 4.2.5 Flood Fill

\vspace{0.25cm}
\sectiontag{Core technique:} Connected components on an implicit grid

\vspace{0.25cm}
\begin{tightitemize}
  \item Cells are nodes in a grid graph.
  \item Edges exist only between orthogonally adjacent cells with the same digit.
  \item Queries ask whether two cells are connected through one digit type.
\end{tightitemize}
\end{frame}

\begin{frame}{How to Recognize the Approach}
\begin{tightitemize}
  \item Running a search for every query can time out.
  \item Precomputing connected component IDs makes every query constant time.
  \item The grid is implicit, so neighbors are generated with direction arrays.
\end{tightitemize}
\end{frame}

\begin{frame}{Algorithm}
\begin{tightitemize}
  \item Read the grid and initialize component[r][c] = -1.
  \item Scan every cell; when it is unlabeled, start DFS with a new component ID.
  \item DFS labels all same-valued cells reachable by four-direction moves.
  \item For each query, convert coordinates from 1-based to 0-based.
  \item Compare component IDs and print binary, decimal, or neither.
\end{tightitemize}
\end{frame}

\begin{frame}{Walkthrough of the Key State}
\begin{tightitemize}
  \item The DFS stack avoids recursion-limit issues on large regions.
  \item The component ID captures both reachability and digit value.
  \item If IDs differ, no valid path exists even if the endpoint digits match.
\end{tightitemize}
\end{frame}

\begin{frame}{Why the Algorithm Is Correct}
\begin{tightitemize}
  \item Flood fill visits exactly the cells reachable from the start through same-valued moves.
  \item Every cell is assigned to exactly one maximal same-valued connected component.
  \item Two queried cells are connected exactly when their component IDs match.
\end{tightitemize}
\end{frame}

\begin{frame}{Complexity and Implementation Notes}
\small
\sectiontag{Complexity:} O(RC + Q) time and O(RC) memory.

\vspace{0.3cm}
\begin{tightitemize}
  \item Boundary checks must happen before indexing the grid.
  \item Print binary for component value 0 and decimal for component value 1.
\end{tightitemize}
\end{frame}



% BEGIN GENERATED CODE WALKTHROUGH
\begin{frame}[fragile]{Code: Read the Grid}
\scriptsize
\sectiontag{Source lines:} \texttt{10kindsofpeople.py:37--59}

\vspace{0.08cm}
\begin{columns}[T,totalwidth=\textwidth]
\begin{column}{0.58\textwidth}
\begin{lstlisting}[style=cpcode]
import sys

def main():
    input = sys.stdin.buffer.readline

    rows, cols = map(int, input().split())

    grid = []

    for _ in range(rows):
        grid.append(input().strip().decode())
\end{lstlisting}
\end{column}
\begin{column}{0.38\textwidth}
\sectiontag{What this chunk does}
\begin{tightitemize}
  \item The first line gives the grid dimensions.
  \item Each row is stored as a string of 0 and 1 characters.
  \item The digit represents which type of traveler can use that cell.
\end{tightitemize}
\end{column}
\end{columns}
\end{frame}

\begin{frame}[fragile]{Code: Prepare Components}
\scriptsize
\sectiontag{Source lines:} \texttt{10kindsofpeople.py:60--80}

\vspace{0.08cm}
\begin{columns}[T,totalwidth=\textwidth]
\begin{column}{0.58\textwidth}
\begin{lstlisting}[style=cpcode]
component = [[-1] * cols for _ in range(rows)]

component_id = 0

directions = [(-1, 0), (1, 0), (0, -1), (0, 1)]  # up  # down  # left  # right
\end{lstlisting}
\end{column}
\begin{column}{0.38\textwidth}
\sectiontag{What this chunk does}
\begin{tightitemize}
  \item component[r][c] records the connected-region ID for each cell.
  \item -1 means the cell has not been assigned yet.
  \item The four directions enforce non-diagonal movement.
\end{tightitemize}
\end{column}
\end{columns}
\end{frame}

\begin{frame}[fragile]{Code: Start DFS for One Region}
\scriptsize
\sectiontag{Source lines:} \texttt{10kindsofpeople.py:82--99}

\vspace{0.08cm}
\begin{columns}[T,totalwidth=\textwidth]
\begin{column}{0.58\textwidth}
\begin{lstlisting}[style=cpcode]
def dfs(start_row, start_col, current_id):
    value = grid[start_row][start_col]

    stack = [(start_row, start_col)]

    component[start_row][start_col] = current_id
\end{lstlisting}
\end{column}
\begin{column}{0.38\textwidth}
\sectiontag{What this chunk does}
\begin{tightitemize}
  \item The DFS captures the digit value of the starting cell.
  \item A stack avoids Python recursion-depth problems on large regions.
  \item The starting cell is marked immediately to avoid duplicates.
\end{tightitemize}
\end{column}
\end{columns}
\end{frame}

\begin{frame}[fragile]{Code: Expand Same-Value Neighbors}
\scriptsize
\sectiontag{Source lines:} \texttt{10kindsofpeople.py:100--130}

\vspace{0.08cm}
\begin{columns}[T,totalwidth=\textwidth]
\begin{column}{0.58\textwidth}
\begin{lstlisting}[style=cpcode]
while stack:
    row, col = stack.pop()

    for dr, dc in directions:
        new_row = row + dr
        new_col = col + dc

        if new_row < 0 or new_row >= rows or new_col < 0 or new_col >= cols:
            continue

        if component[new_row][new_col] != -1:
            continue

        if grid[new_row][new_col] != value:
            continue

        component[new_row][new_col] = current_id
        stack.append((new_row, new_col))
\end{lstlisting}
\end{column}
\begin{column}{0.38\textwidth}
\sectiontag{What this chunk does}
\begin{tightitemize}
  \item Each popped cell checks its four neighbors.
  \item Out-of-bounds, already visited, and different-valued cells are skipped.
  \item Valid neighbors receive the same component ID and are explored later.
\end{tightitemize}
\end{column}
\end{columns}
\end{frame}

\begin{frame}[fragile]{Code: Precompute Every Region}
\scriptsize
\sectiontag{Source lines:} \texttt{10kindsofpeople.py:132--149}

\vspace{0.08cm}
\begin{columns}[T,totalwidth=\textwidth]
\begin{column}{0.58\textwidth}
\begin{lstlisting}[style=cpcode]
for row in range(rows):
    for col in range(cols):

        if component[row][col] == -1:
            dfs(row, col, component_id)

            component_id += 1
\end{lstlisting}
\end{column}
\begin{column}{0.38\textwidth}
\sectiontag{What this chunk does}
\begin{tightitemize}
  \item The double loop visits every cell in the map.
  \item An unvisited cell starts a new connected component.
  \item Each component ID is used later for constant-time query answers.
\end{tightitemize}
\end{column}
\end{columns}
\end{frame}

\begin{frame}[fragile]{Code: Answer Queries}
\scriptsize
\sectiontag{Source lines:} \texttt{10kindsofpeople.py:151--195}

\vspace{0.08cm}
\begin{columns}[T,totalwidth=\textwidth]
\begin{column}{0.58\textwidth}
\begin{lstlisting}[style=cpcode]
queries = int(input())

output = []

for _ in range(queries):
    r1, c1, r2, c2 = map(int, input().split())

    r1 -= 1
    c1 -= 1
    r2 -= 1
    c2 -= 1

    if component[r1][c1] != component[r2][c2]:
        output.append("neither")

    elif grid[r1][c1] == "0":
        output.append("binary")

    else:
        output.append("decimal")

sys.stdout.write("\n".join(output))
\end{lstlisting}
\end{column}
\begin{column}{0.38\textwidth}
\sectiontag{What this chunk does}
\begin{tightitemize}
  \item Kattis coordinates are converted from 1-based to 0-based.
  \item Different component IDs mean neither traveler type can get through.
  \item Matching IDs are labeled binary or decimal based on the starting cell's digit.
\end{tightitemize}
\end{column}
\end{columns}
\end{frame}
% END GENERATED CODE WALKTHROUGH

\begin{frame}{Source}
\small
\sectiontag{Solution file:} \texttt{practice\_contests/day\_4/10kindsofpeople.py}

\vspace{0.3cm}
\sectiontag{Problem:} \url{https://open.kattis.com/problems/10kindsofpeople}

\vspace{0.45cm}
Use this deck with the source code open beside it. The slides explain the reasoning; the code shows the exact input handling and output formatting.
\end{frame}

\end{document}
